\section{Introduction}
\label{sec:intro}

% ---------------------------------------------------------------------------
% ~900 words. Argument order, one paragraph each:
%
%  1. Stochastic IMU calibration fits N, B, K from an Allan or wavelet
%     analysis and transfers them across configuration by a sqrt(ODR) rule.
%  2. That rule is stated as a CONDITION in Kalibr's own documentation --
%     requiring ideal decimation -- not as a validated result.
%  3. The standards construct that WOULD model register quantisation, the
%     IEEE-952 angle-quantisation term Q, describes angle-INCREMENT outputs.
%     Wrong architecture for a rate register.
%  4. So for a rate register the term is unmodelled, lands on the -1/2 family,
%     and is absorbed silently into fitted ARW.
%  5. Contributions, four bullets, in Concept Note v2.3 Z.3's survivability
%     order (strongest first).
%
% WHAT IS NOT CLAIMED -- one sentence each, IN the introduction, not buried:
%   - not a new quantiser theory (Bennett, Sripad-Snyder, Widrow-Kollar)
%   - scale factor / misalignment are a different family (IMU-TK)
%   - the effect vanishes at high rho; this is a low-rho paper
% ---------------------------------------------------------------------------

\gap{Kalibr verbatim quote. Concept Note v2.3 Z.5 names this the paper's
weakest structural point and Critical Analysis v2.1 Z.1 makes it a condition
of publishability. Half a day at a desk. Until it is here, \cref{sec:conseq}
overreaches.}

\todo{Write.}
